기룡이는 오버워치 워크숍에서 ``랜덤 마라톤''을 즐기고 있다. 
랜덤 마라톤은 무작위로 배정받은 캐릭터로 모든 거점을 점령한 뒤 기지로 돌아오는 경주이다. 
이번에 기룡이에게 배정된 캐릭터는 ``리퍼''이다.

경주가 열리는 공간은 장애물이 없는 무한한 2차원 평면이다. 
기지 한 곳과 거점 $N-1$곳이 서로 다른 위치에 있다. 
입력으로 주어지는 $N$개의 위치 중 $1$번은 기지이고, 
$2$번부터 $N$번까지는 거점이다.

이 게임의 시간 단위는 틱이며, 모든 행동은 정수 틱 시각에 시작하고 끝난다.

거점 $i$를 점령하려면 해당 위치에서 점령을 시작한 뒤 $o_i$틱 동안 연속해서 머물러야 한다. 
점령하는 동안에는 이동하거나 스킬을 사용할 수 없다.
$o_i=0$이면 점령을 시작하는 즉시 완료된다. 

기룡이는 다음과 같이 이동할 수 있다.
\begin{itemize}
\item \textbf{걷기:} $1$틱 동안 현재 위치에서 유클리드 거리가 $1$ 이하인 원하는 위치로 이동할 수 있다. 따라서 거리가 $d$인 두 위치 사이를 걷기만 해서 이동하려면 최소 $\lceil d\rceil$틱이 필요하다.
\item \textbf{그림자 밟기:} 현재 위치에서 유클리드 거리가 $M$ 이하인 원하는 위치로 즉시 순간 이동한다. 스킬 사용에는 시간이 걸리지 않지만, 한 번 사용하면 그 시각부터 $K$틱이 지난 뒤 다시 사용할 수 있다.
\end{itemize}

그림자 밟기 전후의 걷기 구간은 각각 틱 단위로 계산한다. 예를 들어 $a$에서 $k$까지 걸은 뒤 그림자 밟기로 $k'$으로 이동하고, 다시 $b$까지 걷는다고 하자. 두 걷기 구간의 유클리드 거리가 각각 $d_1$, $d_2$라면 걷기에만 최소 $\lceil d_1\rceil+\lceil d_2\rceil$틱이 필요하다. 스킬의 쿨타임으로 인해 대기해야 한다면 그 시간도 별도로 소요된다.

그림자 밟기는 기지나 거점뿐 아니라 그 사이의 임의의 위치에서도 사용할 수 있다. 기룡이는 원하는 위치에서 대기할 수 있으며, 걷거나 대기하거나 거점을 점령하는 동안에도 스킬의 쿨타임은 줄어든다. 쿨타임이 끝나는 시각부터 스킬을 다시 사용할 수 있다.

기룡이는 $0$틱에 기지에서 출발하며, 처음에는 그림자 밟기를 즉시 사용할 수 있다. 이동 중에 기지나 이미 점령한 거점을 다시 방문해도 된다. 모든 거점을 한 번씩 점령한 뒤 기지로 돌아오는 순간 경주가 끝난다.

기룡이가 경주를 끝내는 데 필요한 최소 시간을 구해 보자.
